\documentclass[12pt]{article}
\usepackage[utf8]{inputenc}
\usepackage[T1]{fontenc}
\usepackage{hyperref}
\title{Siraya Language Resources for K-12 Curriculum Development}
\author{Compiled by Claude Sonnet 5.5 (claude-sonnet-5-5) for Andr\'as Kornai}
\date{}
\begin{document}
\maketitle

\section{Introduction}
Siraya (ISO 639-3 \texttt{fos}) is an Austronesian language of southwestern Taiwan, last spoken natively around 1908 according to many linguists, though activists prefer not to call it extinct. Revitalization draws on Dutch-era texts, chiefly a translation of the Gospel of Matthew (1661).

\section{Primary Educational Resources}
A 2025 article by Christopher Joby reports that activists have used the Dutch texts to produce textbooks, songbooks and other learning materials; specific titles were not identified in these searches. See the UEA record: \url{https://ueaeprints.uea.ac.uk/id/eprint/99055}.

\section{Linguistic and Reference Materials}
\begin{itemize}
\item Alexander Adelaar, \emph{Siraya: Retrieving the Phonology, Grammar and Lexicon of a Dormant Formosan Language} (De Gruyter, 2011), based on the Siraya Gospel of Matthew: \url{https://mitpressbookstore.mit.edu/book/9783110252958}.
\item A review excerpt noting uncertainty about whether the Gospel text reflects native Siraya or learner features of the Dutch ministers: \url{https://muse.jhu.edu/article/530556}.
\end{itemize}

\section{Religious and Community Resources}
The Siraya Gospel of Matthew is the key historical text (see Adelaar above). Mongabay (2026) describes the community revival, led by Uma Talavan and Edgar Macapili, which helped secure recognition as Taiwan's 17th Indigenous group: \url{https://news.mongabay.com/2026/08/siraya-language-deemed-extinct-over-a-century-ago-finds-its-voice-again-in-taiwan/}. Taiwan Today also covers the topic: \url{https://taiwantoday.tw/AMP/society/top-news/15091/{{url}}}.

\section{Digital and Audio Resources}
No audio or digital learning resources were identified in these searches.

\section{Curriculum Development Context}
There is no evidence here of Siraya in schools. Curriculum work would rest on the Dutch-era texts and Adelaar's reconstruction, with community leadership. Existing activist textbooks should be located and verified.
\end{document}